% Generated draft: requires mathematical and source review.
\documentclass[11pt]{article}
\usepackage[margin=25mm]{geometry}
\usepackage{fontspec}
\setmainfont{Latin Modern Roman}
\usepackage{amsmath,amssymb,mathtools,unicode-math}
\setmathfont{Latin Modern Math}
\usepackage{enumitem,needspace,xurl,hyperref}
\setlength{\parindent}{0pt}
\setlength{\parskip}{5pt}
\setlength{\emergencystretch}{3em}
\newcommand{\N}{\mathbb{N}}
\newcommand{\Z}{\mathbb{Z}}
\newcommand{\Q}{\mathbb{Q}}
\newcommand{\R}{\mathbb{R}}
\newcommand{\C}{\mathbb{C}}
\newcommand{\F}{\mathbb{F}}
\newcommand{\A}{\mathbb{A}}
\newcommand{\PP}{\mathbb{P}}
\newcommand{\E}{\mathbb{E}}
\newcommand{\eps}{\varepsilon}
\newcommand{\dd}{\,\mathrm{d}}
\begin{document}
\section*{UnsolvedMath ID: 30003472}
\textbf{Source problem code:} catalogue entry
\subsection{Definitions and mathematical statement}
Let \(\mu\) be a probability measure on \(\R^2\) such that \(\mu(\ell)=0\) for every line \(\ell\subset\R^2\). For a positive integer \(n\), let \(X_1,\ldots,X_n\) be independent random points with common distribution \(\mu\), and write \(X=(X_1,\ldots,X_n)\). These points are almost surely distinct, with no three collinear.

For a labeled configuration \(x=(x_1,\ldots,x_n)\) in general position, define its orientation data by
\[
\chi_x(i,j,k)=\operatorname{sgn}\det(x_j-x_i,x_k-x_i)
\]
for distinct indices \(i,j,k\). A labeled oriented order type records these signs. An unlabeled oriented order type records them up to a permutation of the indices. A convention that also identifies a configuration with its mirror image additionally identifies \(\chi_x\) with \(-\chi_x\). The supplied record does not specify which convention is intended.

Once an order-type convention is fixed, let \(\mathcal T_n\) be the corresponding collection of realizable planar order types on \(n\) points in general position, and define
\[
p_{\mu,n}(\omega)=\Pr[X\text{ realizes }\omega],\qquad \omega\in\mathcal T_n.
\]
The cleaned source asks for an absolute constant \(c>1\) and order types \(\omega_1,\omega_2\in\mathcal T_n\) satisfying
\[
p_{\mu,n}(\omega_1)>c^n p_{\mu,n}(\omega_2).
\]
Here an absolute constant is understood to be independent of both \(\mu\) and \(n\); the selected order types may depend on \(\mu\) and \(n\).

\textbf{Unresolved quantification.} The source does not quantify \(n\) or give a range or lower threshold. If the intended range were a specified set \(I\subseteq\N\), the uniform formulation would be
\[
\exists c>1\;\forall\mu\text{ charging no lines}\;\forall n\in I\;\exists\omega_1,\omega_2\in\mathcal T_n:
\quad p_{\mu,n}(\omega_1)>c^n p_{\mu,n}(\omega_2).
\]
This is a formulation schema, not an assertion that a particular range \(I\) was supplied. An eventual-in-\(n\) formulation would also require specifying whether its threshold may depend on \(\mu\).

\textbf{Unresolved source qualifications.} The original statement requires only \(c>0\), whereas the cleaned statement requires an absolute \(c>1\); these conditions are not equivalent. No positivity condition on \(p_{\mu,n}(\omega_2)\) is supplied. Consequently, the displayed inequality permits a zero-probability second type and must not silently be replaced by a ratio of positive probabilities. The order-type convention, the quantification of \(n\), and the discrepancy concerning \(c\) require review before a unique precise conjecture can be stated.
\subsection{Short English statement}
For independent points drawn from a planar probability measure assigning zero mass to every line, the cleaned record asks whether two order types can have probabilities separated by a factor exceeding \(c^n\), with a universal \(c>1\). The range of \(n\), the order-type convention, and whether both probabilities must be positive are unspecified. The original record instead says \(c>0\). This is a statement-only draft requiring human mathematical review.
\subsection{Sources}
\begin{itemize}
\item Supplied local UnsolvedMath registry record, permanent ID \(30003472\), title \emph{Exponential Probability Gaps Between Random Order Types}, source problem code \texttt{catalogue entry}, in the \emph{Oberwolfach Reports Open Problems} collection. The proposed year is recorded as 2017; no individual proposer is supplied.
\item Supplied report citation: \emph{Discrete Geometry (2018)}, Oberwolfach Reports, DOI 10.4171/owr/2017/19, pp.\ 1135--1205. \url{https://doi.org/10.4171/owr/2017/19}. The background describes this report as documenting workshop discussions and open problems. Its title year and DOI year are reproduced as supplied.
\item Supplied literature reference: Goaoc and collaborators on limits of order types. \url{https://arxiv.org/abs/1811.02236}.
\item Supplied literature reference: SoCG primary paper. \url{https://doi.org/10.4230/LIPIcs.SOCG.2015.300}.
\item Literature assessment dated exactly \texttt{2026-08-21}: the supplied curation labels the problem open and reports a fixed-size probability bias with a ratio exceeding \(1.8\) for six points, while reporting that the cited literature does not establish the requested exponential comparison for every measure and size. This is supplied dated curation derived from the OWR record, not an independently verified literature assessment; its status and cited claims require independent verification.
\item The registry labels the statement \texttt{corrected\_verified} and describes its verification as repairing notation, typography, and sentence structure while preserving mathematical content. This supplied label does not resolve the recorded change from \(c>0\) to an absolute \(c>1\), and is not treated here as independent mathematical verification.
\end{itemize}
\end{document}
